% Part: proof-theory
% Chapter: proof-search
% Section: introduction

\documentclass[../../../include/open-logic-section]{subfiles}

\begin{document}

\olfileid{pt}{ps}{int}

\olsection{جان پچھان}

!!{proof} نظاماں، جیویں \Log{G3c} یا \Log{N1i}، دے مسلّمے تے استنباط دے قاعدے طے کردے نیں کہ سیکوئنٹاں یا !!{formula}s دے کیہڑے درخت درست !!{proof}s منّے جاندے نیں۔ جے سیکوئنٹاں یا !!{formula}s دا ایہو جیہا درخت دِتّا ہووے (شاید نال ہور جانکاری وی ہووے، جیویں ہر قدم اُتے ورتے قاعدے، یا مفروضیاں تے اوہناں نوں خارج کرن والے استنباطی قاعدیاں دے لیبل)، تاں مسلّمیاں تے قاعدیاں دیاں تعریفاں دی مدد نال اسیں ویکھ سکدے آں کہ دِتّا درخت درست !!{proof} اے کہ نہیں۔ !!a{proof} \emph{لبھنا}، کسے دِتّے سیکوئنٹ دا یا دِتّے مفروضیاں توں دِتّے !!{formula} دا، وکھری گل اے تے عام طور اُتے اوکھی وی۔ ایہو \emph{ثبوت لبھّن} دا مسئلہ اے۔

ثبوت لبھّن دا اک بہت سدھا طریقہ اے۔ اسیں !!{formula}s نوں اک محدود حرف مالا دیاں نشانیاں دیاں لڑیاں سمجھ سکدے آں۔ سیکوئنٹ تے سیکوئنٹاں یا !!{formula}s دے درخت وی !!{formula}s دیاں لڑیاں سمجھے جا سکدے نیں (شاید درخت دیاں شاخاں وکھاؤن لئی خاص قوساں ورت کے)۔ مسلّمیاں تے قاعدیاں دی مدد نال اسیں مشینی طریقے نال جانچ سکدے آں کہ نشانیاں دی لڑی کسے درست !!{proof} نوں ظاہر کردی اے کہ نہیں۔ سو اسیں \emph{سارے ممکن} درست !!{proof}s مشینی طریقے نال بنا سکدے آں (اوہ حرف مالا جیہڑی !!{proof}s ظاہر کرن لئی ورتدے آں، اوہدیاں نشانیاں دیاں ساریاں لڑیاں بنا کے، تے ہر لڑی لئی جانچ کے کہ اوہ کسے درست !!{proof} نوں ظاہر کردی اے کہ نہیں)۔ سدھا طریقہ ایہ اے کہ سارے درست !!{proof}s بناؤندے جاؤ، جدوں تک اک ایہو جیہا !!{proof} نہ مل جاوے جیہڑا دِتّے سیکوئنٹ دا اے (یا دِتّے !!{formula} دا، جیہدے کھُلّے مفروضے دِتّے ہویاں وچوں ہون)۔

چنگا طریقہ اوہ اے جیہڑا ہتھ نال !!a{proof} لبھّن لئی ورتی دا اے۔ آخری سیکوئنٹ توں شروع کرو، !!a{formula} چُنو تے استنباط دا قاعدہ «پچھاں ول» لاؤ، تاں جو اک یا ودھ مقدمے بنن۔ جے اوہناں دے !!{proof}s ہون، تاں اوہ آخری استنباط نال رل کے دِتّے سیکوئنٹ دا !!a{proof} بنا دیندے۔ قدرتی اخذ دے نظام وچ !!a{proof} لبھّن لئی وی رلدا ملدا، پر ذرا اوکھا، طریقہ ورتی دا اے۔

پر ایس طریقے نال کامیابی دی ضمانت نہیں۔ جے غلط قاعدہ چُنو یا !!{formula}s نوں غلط ترتیب نال چُنو، تاں اگلا راہ بند ہو سکدا اے۔ طریقہ پورا متعین وی نہیں۔ ہر قدم اُتے چُنن لئی اک توں ودھ !!{formula} یا پچھاں ول لاؤن لئی اک توں ودھ قاعدے ہو سکدے نیں۔ \emph{ثبوت لبھّن دا الگورتھم} اوہ پورا متعین طریقہ اے جیہڑا !!a{proof} ضرور لبھ لیندا اے جے اوہ موجود ہووے، یعنی اوہدی کامیابی دی ضمانت اے۔

کجھ !!{proof} نظام ہورناں نالوں سوکھے ثبوت لبھّن والے الگورتھم دی گنجائش دیندے نیں۔ سیکوئنٹ نظاماں وچ !!{main operator} جیہڑا !!a{formula} دا اے، تے ایہ گل کہ اوہ کھبّے پاسے اے یا سجّے، طے کردے نیں کہ کِہڑا قاعدہ پچھاں ول لانا اے۔ مثال لئی جے تسی !!a{proof} لبھنا چاہنے او $\Gamma \Sequent \Delta, !A \land !B$ دا، جیہدے وچ $!A \land
!B$ سجّے پاسے اے، تاں \RightR{\land} نوں پچھاں ول لاؤ تے دو میل کھاندے مقدمے $ \Gamma \Sequent \Delta, !A$ تے $ \Gamma \Sequent \Delta, !B$ بناؤ۔ پر جے تہاڈا نظام \Log{G1c} اے، تاں سکڑاؤ ہون نال کم اوکھا ہو جاندا اے کیونکہ چُون لئی ہور راہ ملدے نیں۔ تسی \RightR{\Contraction} نوں پچھاں ول لا کے اوس دی تھاں مقدمہ $ \Gamma \Sequent \Delta, !A \land
!B, !A \land !B$ وی بنا سکدے او۔ فیر مقدمہ $ \Gamma \Sequent
\Delta, !A \land !B, !A \land !B, !A \land !B$، تے ایویں اگے۔ \Log{G2c} وچ ایہ ہور وی اوکھا اے۔ قاعدہ \RightR{\land} پچھاں ول لاؤن دا ایہ وی مطلب اے کہ $ \Gamma_1$، $ \Gamma_2$ دا اندازہ لانا پینا اے، جیہڑے $ \Gamma
= \Gamma_1 \cup \Gamma_2$ پورا کرن، تے $ \Delta_1$، $ \Delta_2$ دا، جیہڑے $ \Delta
= \Delta_1 \cup \Delta_2$ پورا کرن۔ فیر مقدمیّاں $ \Gamma_1 \Sequent \Delta_1, !A$ تے $ \Gamma_2 \Sequent \Delta_2, !B$ نال کوشش کرو۔ غلط اندازے نال اگلا راہ بند ہو سکدا اے تے تہانوں شروع ول مُڑ کے ہور $ \Gamma_1$، $ \Gamma_2$، $ \Delta_1$، $ \Delta_2$ چُننے پین گے۔ جے !!{formula} $!A \lor !B$ اے، تاں \RightR{\lor} دے دو ممکناں وچوں اک چُننا پینا اے، جیہڑا دو ممکن مقدمیّاں وچوں اک $ \Gamma \Sequent \Delta,
!A$ یا $ \Gamma \Sequent \Delta, !B$ بناندا اے۔ غلط چُون نال فیر پچھاں مُڑنا پینا اے۔

\Log{G3c} نظام وچ ایہ اوکڑاں نہیں۔ اوہدے وچ ساختی قاعدے نہیں، تے استنباطی قاعدے دی چُون !!{main operator} تے اوس پاسے نال طے ہوندی اے جیہدے اُتے !!{formula} آؤندا اے۔ ایس لئی \Log{G3c}، \Log{G1c} یا \Log{G2c} نالوں ثبوت لبھّن لئی ودھ راس آؤندا اے۔ کامیاب بھال نال !!a{proof}، \Log{G3c} وچ، ملے گا، تے اوس !!{proof} دا ترجمہ فیر \Log{G1c} یا \Log{G2c} (یا \Log{N1c} وچ وی) ہو سکدا اے۔ سو \Log{G3c} لئی ثبوت لبھّن دا الگورتھم تے \Log{G3c} !!{proof}s دا ہور نظاماں دے !!{proof}s وچ ترجمہ کرن والے الگورتھم ہون، تاں اوہناں نظاماں وچ ثبوت لبھّن دا مسئلہ وی حل ہو جاندا اے۔

سانوں کیویں پتا لگے کہ ثبوت لبھّن والے الگورتھم دی کامیابی دی ضمانت اے؟ ایہدے لئی وکھاؤنا پینا اے کہ جے الگورتھم !!a{proof} نہ لبھ سکے، تاں کوئی !!{proof} موجود نہیں ہو سکدا۔ آخر جے اسیں ماڑا الگورتھم چُنیے، تاں ایہو جیہیاں صورتاں ہو سکدیاں نیں جتھے اوہ ثبوت نہ لبھے بھاویں ثبوت موجود ہووے۔ سدھے الگورتھم نال ایہ اوکڑ نہیں، کیونکہ اوہ سارے ممکن !!{proof}s وچ بھال کردا اے۔ پر ثبوت لبھّن دا الگورتھم کارگر ہونا چاہیدا اے تے گھٹّو گھٹ کم کرنا چاہیدا اے۔

ثبوت لبھّن والے الگورتھم دی کامیابی دی ضمانت وکھاؤن دا مرکزی خیال ایہ اے کہ جے الگورتھم ناکام ہووے، تاں اوہدی ناکامی دے طریقے توں وکھایا جا سکے کہ سیکوئنٹ دا !!a{proof} ہو ای نہیں سکدا۔ ایہ وکھاؤن دا طریقہ اے کہ سیکوئنٹ معتبر نہیں۔ کلاسیکی پہلے درجے دی منطق وچ معتبر سیکوئنٹ $ \Gamma \Sequent \Delta$ اوہ نیں جیہناں لئی ہر !!{structure} وچ گھٹّو گھٹ اک !!{formula}، $\Gamma$ وچلا، جھوٹھا اے یا گھٹّو گھٹ اک !!{formula}،~$\Delta$ وچلا، سچا اے۔ \Log{G3c} \emph{صحت مند} اے، یعنی ایہ صرف معتبر سیکوئنٹ !!{prove}s۔ ایہ !!{proof}s اُتے استقرا نال ثابت ہوندا اے۔ مسلّمے صاف طور اُتے معتبر نیں، تے جے کسے قاعدے دے مقدمے معتبر ہون، تاں اوہدا نتیجہ وی معتبر اے۔ ثبوت لبھّن والے الگورتھم دی کامیابی دی ضمانت وکھاؤن لئی اسیں وکھاؤنے آں کہ جے !!a{proof} لبھّن دی بھال، $\Gamma \Sequent \Delta$ لئی، ناکام ہو جاوے، تاں اسیں !!a{structure} دی تعریف دے سکدے آں جیہدے وچ سارے !!{formula}s، $\Gamma$ والے، سچے نیں تے سارے !!{formula}s،~$\Delta$ والے، جھوٹھے نیں۔ ایس نال ایہ وی ثابت ہو جاندا اے کہ \Log{G3c} \emph{کامل} اے، یعنی جے $ \Gamma \Sequent \Delta$ معتبر اے، تاں اوہدا اک \Log{G3c}-!!{proof} اے۔ کیونکہ ثبوت لبھّن دی ہر ناکام بھال وکھاؤندی اے کہ $ \Gamma \Sequent \Delta$ غیر معتبر اے، معتبر سیکوئنٹ لئی بھال ضرور کامیاب ہونی اے۔

\end{document}
